\chapter{Transferring the Screen~E Actor to WCCI}
\label{ch:screen-e-transfer}

Chapter~\ref{ch:wcci-transfer} ended on two facts about the Screen~E
configuration that point the same way. Retrained from scratch on WCCI it is the
only screened design that produces a usable policy, and it is the only one whose
actor---not merely its encoder---can be moved to another grid at all, because its
head scores a variable-length candidate set with shared weights rather than
indexing a fixed action list.

This chapter takes that portability seriously and asks how much of it can be
realised. It holds the architecture family fixed and varies what crosses the
grid boundary and what is retrained on the far side.

\section{Experiment design}
\label{sec:wcci-design}

Four settings of the Screen~E actor are varied factorially: the input condition
(raw features and absolute angles, or physical scaling with the angle drop), the
touched-node pooling (joint or type-separated mean), the static action descriptor
(off or on), and the idle action (shared or dedicated scorer). Each of the
sixteen cells runs in two arms.
In the \emph{fine-tuned} arm the encoder and the candidate scorer are both
loaded from the paired \texttt{bus14} run and both keep training; in the
\emph{scratch} arm the same configuration starts from random initialization.
Every other setting is copied from the corresponding \texttt{bus14}
configuration, which is what makes the weights loadable at all.
\\
\\
Transferring the scorer as well as the encoder is what the candidate head makes
possible. Its parameters do not depend on the action count, so one agent's
scorer, \texttt{agent\_0} here, can initialize every agent on the target grid.
\\
\\
Loading needs care in one place. The encoder registers its edge index and node
identifiers as persistent buffers, so they appear in the saved state dictionary
at source-grid shape. They describe the source graph and are supplied per agent
from the target environment's own specification at call time, so they are
dropped during the load, while any other shape disagreement is treated as an
architecture mismatch and raises rather than leaving part of the encoder
randomly initialized.

\begin{table}[H]
  \centering
  \small
  \caption{Transfer study design. The two arms of a pair differ only in the two
  transfer settings.}
  \label{tab:wcci-transfer-design}
  \setlength{\tabcolsep}{5pt}
  \begin{tabular}{L{0.26\textwidth}L{0.66\textwidth}}
    \toprule
    Component & Setting \\
    \midrule
    Environment & \texttt{bus36\_wcci\_nomaint}, difficulty 0, topology actions, 80/20 split with split seed 0 \\
    Action space & Reduced at $k=256$, giving 77 / 256 / 127 / 256 actions per agent \\
    Actor & Shared candidate-scoring head on the disaggregated line-node graph, \texttt{share\_actor\_gnn} \\
    Factors & Input condition $\times$ touched-node pooling $\times$ descriptor $\times$ idle scorer, $2^4=16$ cells \\
    Arms & Fine-tuned encoder and scorer from the paired \texttt{bus14} run, both trainable; identical configuration from scratch \\
    Seeds & 0 \\
    Budget & 15M steps, matching the \texttt{bus14} source runs \\
    Inputs & Both conditions are run: raw features with absolute angles, and physical scaling with \texttt{edge\_diff} \\
    \bottomrule
  \end{tabular}
\end{table}

\noindent Running both input conditions rather than only the corrected one is deliberate.
Comparing the scratch arms with and without the correction isolates whether it
helps or harms on the target grid at all, independently of any transfer effect
and on otherwise identical settings; comparing each scratch arm against its
fine-tuned partner isolates the transfer.

% \section{First transfer results}
% \label{sec:wcci-first-results}

% The runs reported here are the first arm described above: the eight launched
% before the input analysis was complete, one seed each, with physical scaling and
% running standardization both disabled and the original angle representation. They
% are therefore the uncorrected-input condition, which is the one
% Section~\ref{sec:wcci-interpretation} predicts should be hardest on a frozen
% encoder. The corrected runs of Table~\ref{tab:wcci-transfer-design} are still
% training.

% Each value below is a deterministic evaluation of a run's best checkpoint over 50
% held-out WCCI chronics. The reference to beat is the do-nothing policy, which
% survives 56\% of those episodes; on this grid inaction is a far stronger baseline
% than on \texttt{bus14}, where it survives about 20\%.

% \begin{figure}[H]
%   \centering
%   \includegraphics[width=0.92\linewidth]{figures/tr_wcci_full_test_bars.png}
%   \caption{50-episode full-test survival of the best checkpoint, by architecture
%   and training arm, with uncorrected inputs and one seed. The dashed line is the
%   do-nothing policy at 56\%.}
%   \label{fig:wcci-transfer-bars}
% \end{figure}

% \begin{table}[H]
%   \centering
%   \small
%   \caption{50-episode full-test survival and the step at which each run's best
%   checkpoint was saved. One seed per cell, uncorrected inputs.}
%   \label{tab:wcci-transfer-results}
%   \setlength{\tabcolsep}{5pt}
%   \begin{tabular}{lrrrrr}
%     \toprule
%     Architecture & Scratch (\%) & Frozen (\%) & $\Delta$ (pp) &
%     Scratch ckpt (M) & Frozen ckpt (M) \\
%     \midrule
%     \texttt{e0n0v0} (bus) & 47.62 & 30.48 & $-17.13$ & 7.38 & 0.58 \\
%     \texttt{e1n0v0} (bus) & 51.31 & 13.94 & $-37.37$ & 4.31 & 13.02 \\
%     \texttt{gb2a\_la2b} & 52.32 & 52.60 & $+0.27$ & 7.38 & 14.68 \\
%     \texttt{gb2a\_lb2a} & \textbf{54.46} & 54.37 & $-0.09$ & 7.38 & 10.12 \\
%     \midrule
%     Do-nothing policy & \multicolumn{2}{c}{56.0} & & & \\
%     \bottomrule
%   \end{tabular}
% \end{table}

% \begin{figure}[H]
%   \centering
%   \includegraphics[width=\linewidth]{figures/tr_wcci_frozen_vs_scratch_curves.png}
%   \caption{Test episodic survival during training, one panel per architecture,
%   frozen transfer against training from scratch. The dashed line is the
%   do-nothing reference at 57.2\% on the training-time evaluation split.}
%   \label{fig:wcci-transfer-curves}
% \end{figure}

% \noindent\textbf{No arm beats doing nothing.} The best of the eight runs reaches
% 54.5\% against 56.0\% for a policy that never acts, and in
% Figure~\ref{fig:wcci-transfer-curves} the curves cross the reference only in
% transient peaks. Whatever the graph representation contributes on
% \texttt{bus14}, it does not yet produce a useful controller here, and every
% comparison below is between policies that are all worse than inaction.
% \\
% \\
% \textbf{Freezing is neutral on the disaggregated schema and destructive on the
% bus schema.} The two heterogeneous architectures match their scratch
% counterparts to within 0.3 points, so the transferred encoder neither helps nor
% obstructs. The two bus architectures lose 17.1 and 37.4 points. This is the
% deficit Section~\ref{sec:wcci-interpretation} warned would be confounded, and the
% confound is real: with uncorrected inputs a frozen \texttt{bus14} encoder is
% reading displaced power and angle channels, so the loss cannot be attributed to
% the representation failing to generalize. The asymmetry between the schemas is
% the part worth pursuing. In the bus graph every node carries the power and angle
% channels, because generator and load values are aggregated onto their busbar,
% so a frozen encoder meets the miscalibrated quantities on every node it reads.
% In the disaggregated graph those channels are non-zero only on asset nodes,
% which are a minority of the node set, while busbars carry cooldown and
% connectivity, and \texttt{rho} rides on the edges where
% Section~\ref{sec:wcci-interpretation} measured it as already comparable across
% grids. That is a hypothesis this experiment does not test, but it is the one the
% corrected runs will separate.
% \\
% \\
% \textbf{The bus architectures are also the unstable ones.} Their end-of-training
% survival, 10\% and 18\% on the smoothed curve, is far below their own best
% checkpoints at 47.6\% and 51.3\%: both peak early and decay for the rest of
% training. The disaggregated runs end near 47\% and 51\%, close to their best
% checkpoints, so their endpoint numbers describe a policy that still exists at
% the end of training rather than one that was passed through. The gap is worth
% keeping in view whenever a best-checkpoint statistic is quoted for this
% environment.
% \\
% \\
% \textbf{The \texttt{bus14} ordering partly survives.} The two disaggregated
% architectures finish above the two bus ones, and \texttt{gb2a\_lb2a}, the
% direction pair that won the message-direction screen, is the best of the four
% here as well. That is consistent with Chapter~\ref{ch:experiments} without being
% evidence for it: one seed per cell and 2.1 points between the two disaggregated
% architectures cannot separate them.
% \\
% \\
% \textbf{What is still open.} One seed per condition, uncorrected inputs, and an
% encoder frozen for all 15M steps rather than unfrozen after a warm start, which
% tests reuse without adaptation and not the only useful form of transfer. The
% first thing the corrected runs have to establish is not the transfer question at
% all, but whether any arm can clear 56\%.


% \section{What this chapter establishes}
% \label{sec:wcci-summary}

% Independently of how the two arms eventually compare, three findings hold.

% \begin{itemize}
%   \item \textbf{The encoder is portable; its inputs are not.} Parameter shapes
%   cross grids unchanged once the grid-shaped buffers are dropped. What does not
%   cross is the distribution those parameters were calibrated on, and nothing in
%   the training configuration was correcting it.

%   \item \textbf{Per-feature standardization can be actively harmful when a
%   channel is a mixture.} Standardizing a channel that is 85\,\% structural zeros
%   aligns its moments and misplaces its most frequent value, degrading every
%   channel it touches. This is not specific to power grids: it applies to any
%   graph representation in which a node feature exists only for nodes carrying a
%   particular kind of asset.

%   \item \textbf{For gauge-dependent quantities, representation beats
%   normalization.} The single largest mismatch measured here was an artifact of
%   the voltage-angle reference. No choice of scaling addresses it, whereas
%   expressing the angle as a difference across each line removes it by
%   construction and makes the channel dense at the same time.
% \end{itemize}
